\documentclass[11pt,a4paper]{article}
\usepackage[margin=2.2cm]{geometry}
\usepackage{booktabs}
\usepackage{enumitem}
\usepackage{xcolor}
\usepackage{hyperref}
\usepackage{amssymb}
\hypersetup{colorlinks=true,linkcolor=black,urlcolor=blue!50!black}
\setlist{itemsep=2pt,topsep=4pt}
\newcommand{\pos}[1]{\textcolor{green!45!black}{#1}}
\newcommand{\negv}[1]{\textcolor{red!65!black}{#1}}

\title{Forward Test: Request for Approval Before Live Paper Trading}
\author{FYP: Learned Market Making on Kalshi Event Contracts}
\date{14 October 2026}

\begin{document}
\maketitle
\vspace{-1.5em}

\paragraph{Purpose.} Before I officially start live paper trading, I would like your approval on two things: the
\textbf{live paper-trading setup} (Section~\ref{sec:infra}) and the \textbf{baselines} it is compared against
(Section~\ref{sec:baselines}). Sections~\ref{sec:evidence}--\ref{sec:tuning} summarise the evidence behind the
headline model and the decision not to tune it. The specific questions are in Section~\ref{sec:questions}.

\section{Strategy and headline model}
The strategy is a passive market maker on Kalshi's multi-outcome markets. Every hour it considers posting a buy or
sell limit order (10 contracts) on each leg, and a learned model decides which orders to post: an order is posted
only if its predicted net profit after fees exceeds 1\,c per contract. The headline model, \textbf{G2}, is a
gradient-boosted tree model on 14 market features (price, spread, recent trades, basket structure, fees,
category), retrained on all data up to August 2026. Its predecessor \textbf{G} (trained on Nov 2024--May 2025)
passed a pre-registered test on Aug--Nov 2025 (Section~\ref{sec:evidence}). All rules, models and statistics were
fixed and timestamped publicly before any live result.

\section{Live paper-trading setup}\label{sec:infra}
\begin{enumerate}[leftmargin=*]
  \item \textbf{Live decisions.} A producer on a Lumid sandbox scores every candidate order each hour from Kalshi's
    live data and records every strategy's decision before outcomes are known. Each hour's decisions are posted to
    Lumid Research Fleet with a SHA-256 hash, giving an independent timestamp.
  \item \textbf{Live fills.} Right after each hour closes, a paper trader determines that hour's fills
    \emph{only from Kalshi trades that have already happened}, using the same fill rules as the backtest:
    \begin{itemize}
      \item an order joining the best price fills only if real trades go \emph{through} its price (others are ahead
        in the queue); an order that improves the price by 1\,c fills on trades at its price;
      \item position limits (50 contracts per leg, 200 per event), Kalshi maker fees, full collateral;
      \item positions are held to settlement and valued at the actual results.
    \end{itemize}
  \item \textbf{Append-only ledger.} Each hour's fills are written to a ledger entry that is never revised. Entries
    are hash-chained (each contains the previous entry's hash). Every round re-checks the recorded fills against a
    recomputation; any mismatch is logged, never overwritten. Settlement profit is added as markets resolve.
  \item \textbf{Hourly reporting.} Each hour, for every strategy, Research Fleet receives the
    \emph{cumulative} realised P\&L (settled markets), unrealised P\&L (open positions at the current mid),
    total P\&L, capital committed and return, plus the ledger entry's hash. Hourly P\&L is the change in total P\&L.
  \item \textbf{Fill assumptions.} Pessimistic fills (above) decide the verdicts; optimistic fills (trades at our
    price also count) and strict fills (competitors match our price improvements) are reported as bounds.
  \item \textbf{What counts.} Only hours recorded live, from the first ledger entry onwards, count towards the
    forward-test result. Earlier hours (11--14 October) are reported separately as a replay. The interim cut for the
    report is 28 October; quoting ends 11 November.
  \item \textbf{Verdicts (pre-registered, one-sided 5\%).} (a) G2 is profitable live: the lower 95\% bound of its net
    return is above zero; (b) G2 adds value: G2 minus random selection at the same volume is above zero. Confidence
    intervals are bootstrapped by event (and by day, as a sensitivity check).
\end{enumerate}
\noindent\textbf{Checks done.} The paper trader reproduced exactly the 1.1 million orders the producer logged
(0 mismatches); two blind test rounds in the production sandbox (no P\&L viewed) confirmed the ledger, the hash
chain, re-verification and an unchanged producer. Some hours on 13 October were lost to an infrastructure
incident; they are recorded as gaps and not repaired.

\section{Baselines}\label{sec:baselines}
All baselines run on the same hours, candidates, fill rules and statistics as G2.
\begin{center}
\small
\begin{tabular}{@{}p{3.0cm}p{5.5cm}p{7.0cm}@{}}
\toprule
Baseline & Rule & Question it answers \\
\midrule
Random selection & random orders at G2's posting volume & Is profit from choosing well, or from trading less? \\
Quote everything & every candidate order & Does any selection help at all? \\
Wide spreads & post if spread $\ge$ 3\,c & Is G2 just a spread filter? \\
Liquid markets & post if recent trades $\ge$ hourly median & Is G2 just a liquidity filter? \\
Linear model & ridge regression, same data and features & Is a nonlinear model needed? \\
Neural network & 2-layer MLP, same data and features & Does another model family do better? \\
Previous model (G) & 2025 model, unchanged & Has the edge decayed? \\
S\&P 500, 3-month T-bill & index and risk-free rate & Better than passive investment? \\
\bottomrule
\end{tabular}
\end{center}
Secondary variants (reported, never verdicts): G2 trained on shorter windows, a two-price-level version, and a
weekly category-selection variant.

\section{Evidence before the live test}\label{sec:evidence}
\paragraph{2025 test (Aug--Nov 2025, pre-registered; full simulator, pessimistic fills).}
\begin{center}
\small
\begin{tabular}{@{}lrrrr@{}}
\toprule
Strategy & Return & 95\% CI & Fills & Net P\&L \\
\midrule
\textbf{G (headline 2025)} & \pos{\textbf{+3.55\%}} & [+2.51, +4.60] & 34,086 & \pos{+\$3,748} \\
Linear model & \pos{+4.03\%} & [+2.94, +5.14] & 18,305 & \pos{+\$2,427} \\
Neural network & \pos{+2.11\%} & [+1.30, +2.92] & 43,975 & \pos{+\$3,104} \\
Wide spreads & +0.07\% & [$-$0.45, +0.55] & 97,517 & +\$206 \\
Liquid markets & $-$0.05\% & [$-$0.51, +0.39] & 99,930 & \negv{$-$\$145} \\
Random (5 seeds, mean) & \negv{$-$0.23\%} & each includes 0 & $\approx$24,400 & \\
Quote everything & \negv{$-$0.34\%} & [$-$0.77, +0.10] & 110,867 & \negv{$-$\$1,133} \\
\bottomrule
\end{tabular}
\end{center}
G minus each baseline (paired by event): random +3.2 to +4.4 (all five 95\% CIs exclude 0); wide spreads +3.47
[+2.50, +4.45]; liquid +3.60 [+2.64, +4.55]; neural +1.43 [+0.47, +2.38]; linear $-$0.48 [$-$1.73, +0.81]
(no significant difference). Under strict fills, G's edge falls to +0.39\% [$-$0.96, +1.78].

\paragraph{2026 walk-forward (Apr--Sep 2026, each month trained only on earlier months; offline scoring).}
G2 made money in all six months (mean +2.65\% per month, +\$3,291), while random selection ($-$1.15\%), both
rules and quoting everything lost money in all six. The linear model averaged +3.76\% (+\$1,455); the neural network
+1.00\%, positive in only three months. In the full simulator on September 2026, G2 made +1.41\% [$-$1.47, +4.31].

\section{Why G2 stays the headline model (linear model)}
The linear model shows a higher return on capital because it trades less: in the 2025 test it posted 13.5\% of
candidate orders against 19.7\% for G, and filled about half as many contracts. G earns more total profit (2025: +\$3,748 vs +\$2,427; 2026: more profit
in every month), and their returns do not differ significantly. When both models post the \emph{same} number of
orders, G2 earns more than both the linear model and the neural network at every volume tested in 2026. In 2025, G
leads at the volume it actually trades (about 20\% of candidates: +5.30\% vs +4.59\% linear and +4.36\% neural);
the models converge at very high volumes, which I interpret as the easier 2025 market rather than overlapping
picks (overlap rises similarly in both years). G2 is therefore kept as the headline model.

\section{Hyperparameters: fixed, not tuned}\label{sec:tuning}
Model settings were fixed in advance and checked rather than tuned:
\begin{itemize}[leftmargin=*]
  \item \textbf{Boosted model:} a walk-forward search over 14 alternative settings placed the default within 0.2
    percentage points of the best (+2.71\% vs +2.91\% per month), below the month-to-month variation. Choosing the
    best of many settings on noisy financial data would inflate its apparent advantage, so the default was kept.
  \item \textbf{Neural network:} across 8 settings, the best averaged +1.45\% against G2's +2.65\%, and every
    setting had losing months; tuning could not change the comparison.
  \item \textbf{Training window:} chosen by a six-month walk-forward (all history +2.65\% vs 3-month +1.53\% and
    6-month +1.88\%).
\end{itemize}
Settings were varied one at a time; a joint search over combinations is listed as future work.

\section{Questions for your approval}\label{sec:questions}
\begin{enumerate}[leftmargin=*]
  \item[$\square$] Does the live paper-trading setup (hourly recording of fills from trades that have already
    happened, append-only ledger, hourly Research Fleet reports) meet the live paper-trading requirement?
  \item[$\square$] Are the baselines in Section~\ref{sec:baselines} sufficient?
  \item[$\square$] Is it acceptable to keep the fixed model settings, given the sensitivity checks in
    Section~\ref{sec:tuning}?
  \item[$\square$] Is G2 as the headline model acceptable, given that the linear model has a higher return but
    lower total profit?
  \item[$\square$] Is anything else you would like reported each hour or at the cuts (28 October, 11 November)?
\end{enumerate}
If you approve, I will start the live ledger immediately; every hour from that point counts towards the result.

\end{document}
